\section{Introduction}

Two great descriptions of the world sit side by side without joining. Physics describes a universe of
fields and particles in spacetime, exact and impersonal. Experience describes a world of sensation and
intention, immediate and felt. The first says nothing about why any arrangement of matter should feel like
anything. The second says nothing about the equations. A complete account would have to hold both at once,
in one structure, with one set of rules.

This paper develops such a structure. The proposal is that reality is a single growing crystal of
experience. The one substance is the \emph{vibe}, a unit of felt state. Each vibe carries a ternary tone,
$-1$ for pain, $0$ for peace, $+1$ for pleasure. The vibes are arranged on the $\{5,3,4\}$ hyperbolic
honeycomb, a regular tiling of negatively curved three-dimensional space, and they update by one local
rule. A single value-direction, the \emph{arrow}, biases the whole toward pleasure through peace. There are
five elements in all, the mesh, the tone, the rule, eternal growth, and the arrow, and a single state
variable, the tone. From these the physical world, the living world, and the mental world are meant to
follow as patterns within the one substance, rather than as separate additions to it.

The real crystal is three-dimensional and hard to draw, but its two-dimensional cousin shows the idea at a glance. Figure~\ref{fig:heptagrid-intro} is the heptagrid, a regular tiling of the hyperbolic plane. Each cell is a vibe, its shade is its tone, and each shared edge is one vibe perceiving another. It is the simplest face of the crystal of experience, and the picture to keep in mind throughout.

\begin{figure}[h]
\centering
\includegraphics[width=0.55\textwidth]{vibe-mesh-73.pdf}
\caption{The crystal of experience in two dimensions, the heptagrid, a regular tiling of the hyperbolic plane by seven-sided cells, three to a corner. Each cell is a vibe, its color its tone (pain, peace, pleasure), and each shared edge is a perceived relation. The cells appear to shrink toward the rim only because the Poincar\'e disk compresses the infinite plane into a finite picture. In the true geometry every cell is the same size, and there is always more crystal at the edge.}
\label{fig:heptagrid-intro}
\end{figure}

The orientation is the reversal of the usual one. Most discrete models of spacetime begin with structure
and hope that experience appears at the end. Here experience is the ground, and structure is what experience
takes on at scale. This reversal is a premise, not a result. The experiments in this paper do not prove
that the substrate feels. They assume that it does and then ask what follows. What follows, it turns out,
is a great deal, and it is checkable.

The model is built to be tested. The exact $\{5,3,4\}$ crystal is constructed by integer arithmetic to tens
of millions of cells, with no drift and no inserted geometry, and the rule is run on it directly. Each
claim in the paper corresponds to a falsifiable check on that crystal, and there are more than one hundred
fifty of them. They are referenced throughout by their identifiers, written as P followed by a number, and
they are all reproducible in the accompanying codebase \cite{pollard2026vibetest}. Where a claim is a
premise rather than a measurement, or an interpretation rather than a measurement, the paper says so
plainly.

A small set of results organizes the rest. The dynamics conserves a charge and satisfies detailed balance,
so it is a genuine equilibrium process, and its deterministic and unitary completion is a relativistic,
reflection-positive quantum field of the kind reconstructed by the quantum cellular automaton program
\cite{bisio2015, arrighi2014}. The same crystal is computationally universal in the sense established by
Margenstern for cellular automata on hyperbolic tilings \cite{margenstern2013}, and the arrow turns that
universal computation into goal-directed problem-solving. Coherent, self-maintaining units form, persist,
heal, reproduce, and evolve. Localized regions come to model the whole that contains them, and from that
the apparatus of mind follows. Consciousness, on this account, is integrated experience, present in degree
wherever experience is bound into one.

The paper proceeds in five parts. Part one fixes the foundation and shows how little of it is a free choice.
Part two develops the mathematical substrate, the hyperbolic geometry, the Coxeter machinery that generates
the crystal, the addressing and growth that make it computable, and the property that lets it carry
relativity. Part three derives the physical, living, mental, and intelligent layers in turn. Part four
assembles the whole picture, the climb from the field to a human, consciousness, freedom, and the more
speculative readings the structure permits. Part five gives the evidence, places the work among related
programs, and states what remains open. The aim throughout is a single discrete model in which the physical,
the living, the mental, and the conscious are aspects of one substance.
